\section{Preliminaries}\label{sec:preliminares}

\paragraph{Board and lines.}
The squares of the board are the pairs $(i,j)$ with $1\le i,j\le n$; $i$ is the row and $j$ the column.
There are four families of lines: the rows $R_i$, the columns $C_j$, the diagonals
$D_k=\{i-j=k\}$, $-n<k<n$, and the antidiagonals $A_l=\{i+j-1=l\}$, $0<l<2n$. The families have $n$, $n$,
$2n-1$ and $2n-1$ lines, respectively, for a total of $6n-2$. Each square lies on exactly four lines, one from each
family, and two lines from different families meet in at most one square.
We number rows and columns from one rather than from zero as in~\cite{Perfil}
and~\cite{Completabilidad}: our square $(i,j)$ is their square $(i-1,j-1)$, our diagonal
$D_k$ has difference $k$ in their notation, and our antidiagonal $A_l$ has sum $l-1$.
Figure~\ref{fig:incidencia} shows how to read all four labels for a single square.
\begin{figure}[htbp]
\centering
\begin{tikzpicture}[x=6.9mm,y=6.9mm]
 \bcuadricula{5}
 \draw[btinta,thin] (0,3.5)--(5,3.5);
 \draw[btinta,thin] (2.5,0)--(2.5,5);
 \draw[btinta,line width=1.2pt] (1,5)--(5,1);
 \draw[btinta,densely dashed,line width=1.1pt] (0,1)--(4,5);
 \fill[bpapel] (2.5,3.5) circle[radius=.26];
 \draw[btinta,line width=.85pt] (2.5,3.5) circle[radius=.24];
 \fill[btinta] (2.5,3.5) circle[radius=.075];
 \node[anchor=west,font=\small] at (6.1,4.25) {Marked square: $(2,3)$};
 \draw[btinta] (6.2,3.35)--(7.1,3.35);
 \node[anchor=west,font=\small] at (7.4,3.35) {Row $R_2$ and column $C_3$};
 \draw[btinta,line width=1.2pt] (6.2,2.35)--(7.1,2.35);
 \node[anchor=west,font=\small] at (7.4,2.35) {Diagonal $D_{-1}$};
 \draw[btinta,densely dashed,line width=1.1pt] (6.2,1.35)--(7.1,1.35);
 \node[anchor=west,font=\small] at (7.4,1.35) {Antidiagonal $A_4$};
\end{tikzpicture}
\caption{The four lines through $(2,3)$. The diagonal and antidiagonal indices
are determined by the coordinates of the square: $2-3=-1$ and $2+3-1=4$. In the incidence matrix,
the row for this square has a one in the column for each of the four lines shown.}
\label{fig:incidencia}
\end{figure}


We write $A$ for the incidence matrix, with one row per square and one column per line.
For $p\in\mathbb R^{n^2}$ and a line $L$, set
\[
 m_L(p)=\sum_{c\in L}p_c,\qquad h_L(p)=1-m_L(p).
\]

The quantity $m_L$ is the mass assigned to $L$, and the amount $h_L$ by which it falls short of one is its vacancy.
For example, if we assign $1/5$ to every square of the board in the figure, each row sums to one,
whereas diagonal $D_{-1}$ sums to $4/5$: its vacancy is $1/5$.

\paragraph{The continuous profile of mass one.}
To compare boards of different sizes, we divide each coordinate by $n$ and work on a fixed square.
The discrete lines then give projections of the profile.
We follow the notation of~\cite{Perfil}. On the square $[0,1]^2$ with coordinates $(r,c)$, the
projections for the four families are $\ell_R=r$, $\ell_C=c$, $\ell_D=r-c$ and $\ell_A=r+c$, with
values in $I_R=I_C=[0,1]$, $I_D=[-1,1]$ and $I_A=[0,2]$. For a density $\nu\ge0$, $m_f$ is its
projection onto family $f$ and $u_f=1-m_f$ its vacancy. The admissible class at mass one,
$\mathcal A_1$, consists of densities of mass one with $m_f\le1$ almost everywhere, and
\[
 s[\nu]=-3-\int_{[0,1]^2}\nu\log\nu-\sum_{f}\int_{I_f}u_f\log u_f
\]
is the functional of~\cite[(1.1)]{Perfil} at mass one. By~\cite[Theorem~1.1]{Perfil}, $s$ attains its
supremum on $\mathcal A_1$ at a unique density $\nu=\nu_1$. We define the optimal value at mass
one, written in~\cite{Perfil} as $s^*(1)$ and here, as in~\cite{Completabilidad}, as $s_1$, and two
associated quantities:
\begin{equation}\label{eq:s1}
 s_1=s[\nu_1]=\sup_{\nu\in\mathcal A_1}s[\nu],\qquad \Phi=s_1+1,\qquad I_2=\int_{[0,1]^2}\nu_1^2 .
\end{equation}
At mass one the row and column vacancies vanish, so only the diagonal vacancies contribute.
According to~\cite[Remark~2.3]{Perfil}, $s$ then has, on admissible densities, the same expression as the entropy
in Simkin's variational problem~\cite[Definition~2.9]{Simkin}. In~\cite{Simkin}, $-\alpha$ is the maximum of that entropy over
a class of measures defined by closure (Theorem~2.24 and Lemma~2.22); Remark~2.7 of that paper specifies,
for absolutely continuous measures, the relation between this class and the class defined by the projection
constraints. The correspondence above and that remark would give $\alpha=-s_1$, that is, $\Phi=1-\alpha$ and $\log Q(n)=\log n!+n\Phi+o(n)$. We call this
the identification with Simkin's problem. We do not use it in any proof: the definition
$\Phi=s_1+1$ and the theorems on the predictor are independent of it.

By~\cite[Theorems~1.2 and~1.3 and Proposition~5.2]{Perfil}, the profile of mass one factors as
\begin{equation}\label{eq:factorizacion}
 \nu_1(r,c)=b_0\,r_R(r)\,r_C(c)\,u_D(r-c)\,u_A(r+c),
\end{equation}
with $b_0>0$ and continuous, positive Lipschitz functions $r_R,r_C,u_D,u_A$, $\int_0^1r_R=\int_0^1r_C=1$.
In those results, these factors are the limits of the rescaled profile at the full board and
are written $b_0$, $r_{R,0}$, $r_{C,0}$, $u_{D,0}$ and $u_{A,0}$; as in~\cite{Completabilidad},
we omit the subscript except in the constant $b_0$.
The row and column marginals of $\nu_1$ equal one, and $u_f+m_f=1$ for $f=D,A$, where
$m_f$ is the projection of $\nu_1$ onto family $f$. There are positive constants
$c_\nu\le C_\nu$, $\kappa_\nu\le1$ and $L_\nu$ such that
\[
 c_\nu\le\nu_1\le C_\nu,\qquad \kappa_\nu\le u_D,u_A\le1,
\]
and $L_\nu$ is a Lipschitz constant for $\nu_1$ with respect to the maximum metric. Sections~\ref{sec:perfil}
and~\ref{sec:constante} use only these properties of the continuous profile. In those sections we write $\nu$
for $\nu_1$.

\paragraph{Conventions.}
The constants $C,c,M,\eta,\kappa,\dots$ do not depend on $n$ unless stated otherwise;
$O(\cdot)$ and $o(\cdot)$ refer to $n\to\infty$. In Sections~\ref{sec:insercion}
to~\ref{sec:meso} the parameter is the number $k$ of empty rows: there $O(\cdot)$ refers to
$k\to\infty$, and each section specifies what the constants depend on. The letter $h$, which
denotes that number in~\cite{Completabilidad}, is used here for vacancies. In Sections~\ref{sec:perfil}
and~\ref{sec:constante} the equations have labels specific to each subsection, such as (P1.3) or (D2.5).
In the normalization logarithms $\log(4\pi n)$ and $\log(2\pi)$, in Stirling's formula, in the
Slater point of Section~\ref{sec:P01} and in Appendix~\ref{sec:torres}, $\pi$ denotes
$3.14159\ldots$; elsewhere, $\pi=\pi^{(n)}$ is the canonical maximizer and $\pi_c$ the weight of
square $c$. Some letters that have a fixed meaning in~\cite{Perfil}, such as $\rho$, $x$, $\lambda$,
$\ell$, $u$, $F$, $J$, $L$ or $\mathcal H$, denote other objects in some sections of this paper.
We define them where they occur; for example, $\rho(B)$ is not a mass.
